\documentclass[a4paper,10pt]{report}

\def\packagepath{../../preambule}   % path du package princial
\usepackage{\packagepath/preambule}    % utilisation du fichier de configuration

\def\level{Première }              % Classe
\def\course{NSI}              % Matière
\def\eval{Complexité -- Tris}

\def\visibleornot{visible}    % visible or invisible
\def\documentpath{./Sources_Latex}
\def\date{Vendredi 28 Mars 2025}
\renewcommand{\arraystretch}{1}  % Ecart dans les tableau

\def\ptsexoA{4}
\def\ptsexoB{4}
\def\ptsexoC{3}
\def\ptsexoD{5}
\def\ptstotal{\ptsexoA+\ptsexoB}

\usepackage{hyperref}

\usetikzlibrary{shapes, arrows, positioning}
\usetikzlibrary{shapes.geometric, arrows, positioning, decorations.pathreplacing}

\tikzstyle{etat} = [draw, rounded corners=15pt, minimum width=2.5cm, minimum height=1.2cm, text centered, font=\bfseries]
\tikzstyle{nouveau} = [etat, fill=blue!40]
\tikzstyle{pret} = [etat, fill=green!60]
\tikzstyle{elu} = [etat, fill=yellow!60]
\tikzstyle{bloque} = [etat, fill=red!60]
\tikzstyle{termine} = [etat, fill=white]
\tikzstyle{fleche} = [thick, ->, >=stealth]

\usepackage{float} % Permet d'utiliser [H]
\usepackage[T1]{fontenc}
\begin{document}

%%%%%%%%%%%%%%%%%%%%%%%%%%%%%%%%%%%%%%%%%%%%%%%%%%ù

%\PageGardeSujetBac[Matiere = Numérique et Sciences Informatiques,
%                  Session = 2025,
%                  AffJour=false,
%                  Duree = {3 heures 30},
%                  DernierePage = \pageref{LastPage},
%                  ModeExamen=false
%                  ]
\renewcommand{\labelitemi}{\textbullet} %pour éviter les tirets dans les "itemize" qui apportent confusion avec le signe moins.
%% Début page de garde style BAC

   %%%%%%%%%%%%%%%%%%%%%%%%%%%%%%%%%%%%%%%%%%%%%%%%%%%%%%%%
%\PageGardeSujetBac[clés]

\pagestyle{DS_FP}          % Feuille de style fancy
\NomPrenomNote{}


\exods{\ptsexoA}

\begin{enumerate}
    \item Ecrire un programme python qui vérifie si une liste \verb|L| de longueur \verb|n| est triée dans l'ordre croissant. Renvoie True si la liste est triée et False sinon.
    On utilisera obligatoirement une boucle. 
    \begin{answer}{1}{8.6}
        \begin{verbatim}
def est_triee(L:list)->bool:
    ''' Renvoie True si L est triée en ordre croissant et False sinon'''
    
    
               
    

    


          
    



    
    
assert est_triee([1, 2, 3, 4, 5]) == True
assert est_triee([1, 2, 8, 4, 5, 6]) == False

        \end{verbatim}
        \end{answer}

        \begin{minipage}{0.49\linewidth}
            \item Donner un exemple de liste de 4 élements pour lequel \verb|est_triee| utilise le plus de boucles.

    \end{minipage}\hfill
    \begin{minipage}{0.49\linewidth}
            \begin{answer}{1}{1}    \end{answer}

    \end{minipage}

    \begin{minipage}{0.49\linewidth}
        \item Donner un exemple de liste de 4 élements pour lequel \verb|est_triee| utilise le moins de boucles.

\end{minipage}\hfill
\begin{minipage}{0.49\linewidth}
        \begin{answer}{1}{1}    \end{answer}

\end{minipage}

\end{enumerate}

\exods{\ptsexoB}

    \begin{enumerate}
\begin{minipage}{0.48\linewidth}
    On considère l'algorithme suivant:

\begin{CodePythontexAlt}[Largeur=1\linewidth,Centre,PremLigne=1]{}
def mystere(mot:str)->bool:
    ''' mot: str
    la valeur renvoyée: bool '''
    i, n = 0, len(mot)
    j = n - 1
    p = True
    while i <= j:
        if mot[i] != mot[j]:
            p = False
        i = i + 1
        j = j - 1
    return p
\end{CodePythontexAlt}
\item Exécuter cet algorithme pour les mots \verb+arbre+ et \verb+radar+. Quelle est la fonction de cet algorithme?
    
    \begin{answer}{1}{3}    \end{answer}

\end{minipage}\hfill
\begin{minipage}{0.48\linewidth}
    
    
    \item Déterminer le variant de boucle de cet algorithme
    
    
    \begin{answer}{1}{1}    \end{answer}
    \item Montrer que cet algorithme se termine
    \begin{answer}{1}{6.5}    \end{answer}
\end{minipage}

\end{enumerate}


      \newpage
  
      \thispagestyle{DS_LP}          % Feuille de style fancy

\exods{\ptsexoC}

\begin{enumerate}

    
\begin{minipage}{0.29\linewidth}
    Fonction qui calcule la valeur de $2^n$

\begin{CodePythontexAlt}[Largeur=1\linewidth,Centre,PremLigne=1]{}
def puissance2(n):
    p = 1
    i = 0
    while i < n:
        p = p * 2
        i = i + 1
    return p
\end{CodePythontexAlt}

\end{minipage}\hfill
\begin{minipage}{0.66\linewidth}
    \item Effectuer la trace de cet algorithme pour \verb|n = 4|.
    
    \medskip

    \begin{tabular}{|c|c|c|c|c|}
        \hline
        p debut boucle& i debut boucle & while &p fin boucle & i fin de boucle\\ \hline
        1&0&True&&\\\hline
        &&&&\\\hline
        &&&&\\\hline
        &&&&\\\hline
        &&&&\\\hline
        &&&&\\\hline


    \end{tabular}
\end{minipage}


    
    
    \item Proposer un invariant de boucle et démontrer que cet algorithme est correct.

    \begin{answer}{1}{3}    \end{answer}

\end{enumerate}



\exods{\ptsexoD}


On considère le programme en pseudo code suivant qui permet de réaliser un tri par insertion en place dans une liste.

\begin{verbatim}
Le premier élement de la liste est trié
pour tous les éléments de la liste non triée:
    k est l'indice de la clé (premier élément de la liste non triée)
    Tant que le bord gauche n'est pas atteint et que la clé est plus petite que l'élément à sa gauche:
        echanger la clé avec l'élément à sa gauche
        l'indice de la clé diminue de 1
\end{verbatim}

\begin{enumerate}
    \item Traduire en langage python cet algorithme.
    

    \begin{answer}{1}{6}
        \begin{verbatim}
def tri_insertion(liste):
    ''' Tri la liste par insertion - Ne renvoie rien '''

        \end{verbatim}
    \end{answer}


    \begin{minipage}{0.48\linewidth}
        \item Dans le tableau ci-dessous montrer le contenu de la liste \verb|[7, 2, 9, 1, 5, 3, 8, 4, 6]| au fur et à mesure que se déroule le tri par insertion.
    
    \begin{tabular}{|c|c|c|c|c|c|c|c|c|c|}
        \hline
        Indice: & 0 & 1 & 2 & 3 & 4 & 5 & 6 & 7 & 8 \\ \hline
        Départ&7&2&9&1&5&3&8&4&6 \\ \hline
        Etape 1&&&&&&&&& \\ \hline
        Etape 2&&&&&&&&& \\ \hline
        Etape 3&&&&&&&&& \\ \hline
        Etape 4&&&&&&&&& \\ \hline
        Etape 5&&&&&&&&& \\ \hline
        Etape 6&&&&&&&&& \\ \hline
        Etape 7&&&&&&&&& \\ \hline
        Etape 8&&&&&&&&& \\ \hline
    \end{tabular}
    \end{minipage}\hfill
    \begin{minipage}{0.48\linewidth}
        \item Montrer que cet algorithme se termine.
    \begin{answer}{1}{6}
        
    \end{answer}
    \end{minipage}
    
\end{enumerate}
%\appelprof{}
   %\include{\documentpath/LP13}
%\rule{0.95\linewidth}{3pt}

%\renewcommand{\thefigure}{1.\arabic{figure}}
%\renewcommand{\thetable}{1.\arabic{table}}
%\input{./Sources_Latex/exo1_ssrep}\newpage
%\setcounter{figure}{0}
%\setcounter{table}{0}



\end{document}








